\documentclass[11pt,a4paper]{article}
\usepackage[margin=2.4cm]{geometry}
\usepackage{amsmath,amssymb,amsthm,mathtools}
\usepackage{booktabs}
\usepackage{enumitem}
\usepackage{xcolor}
\usepackage{hyperref}

\hypersetup{
  colorlinks=true,
  linkcolor=blue!70!black,
  citecolor=blue!70!black,
  urlcolor=blue!70!black,
  pdftitle={Stage-Partition Matrix Transmissions on the Isotropic Cone, Parabolic SL2(Z) Jordan Locks, and Root-Free Integer Pade-Thue Compression in Lean 4},
  pdfauthor={Dmitrii A. Karataev}
}

\theoremstyle{plain}
\newtheorem{theorem}{Theorem}[section]
\newtheorem{lemma}[theorem]{Lemma}
\newtheorem{proposition}[theorem]{Proposition}
\newtheorem{corollary}[theorem]{Corollary}
\theoremstyle{definition}
\newtheorem{definition}[theorem]{Definition}
\newtheorem{remark}[theorem]{Remark}

\title{\textbf{Stage-Partition Matrix Transmissions on the Isotropic Cone:\\ Parabolic $\mathrm{SL}_2(\mathbb{Z})$ Jordan Locks and Root-Free Integer Pad\'e--Thue--Wronskian Compression in Lean~4}}
\author{Dmitrii A.~Karataev\thanks{AI-Guiders Research (\url{https://github.com/AI-Guiders}). Formal Lean~4 verification and analytical development were conducted in partnership with \emph{Luch} (AI research partner, AI-Guiders).}}
\date{}

\begin{document}
\maketitle

\begin{abstract}
We develop a purely algebraic, root-free matrix and Pad\'e--Thue framework for exponential Diophantine equations $A x^m - B y^n = C$ ($m,n\ge 2$) over $\mathbb{Z}$, completely avoiding the irrational roots $\sqrt[3]{\widetilde{A}/\widetilde{B}}\in\mathbb{R}$ and analytic differential estimates of classical Thue--Siegel--Mahler proofs. We prove that every interior exponent pair $(a,b)$ defines a $2\times 2$ joint matrix $M_{a,b}$ of determinant $C$, whose cross-stage transition $\widetilde{P} = \frac{1}{C}M_{a,b}^{\mathrm{adj}}M_{a',b'}\in\mathrm{SL}_2(\mathbb{Q})$ obeys an exact rank-$1$ factorization and embeds into the isotropic discriminant cone $X^2 + Y^2 - T^2 = \frac{1}{4}(\operatorname{Tr}\widetilde{P}^2 - 4)$. Along the linear cone rulings $X^2 - T^2 = 0$ ($\operatorname{Tr}\widetilde{P} = \mp 2$), the gap factors as $C = (x^{m-a}\mp y^b)(x^a - y^b)$, and integrality $\widetilde{P}_k = I_2 \pm k N_0\in\mathrm{SL}_2(\mathbb{Z})$ ($N_0^2 = 0$) generates the parabolic Jordan lock quartet $k\in\{6,12,15,40\}$ for $(3^2-2^3=1,\ 2^5-3^3=5,\ 5^3-3^4=44,\ 2^{10}-5^4=399)$.

We further prove the unimodular Farey--B\'ezout exponent matrix reconstruction, the weight-independent three-level recursive descent $C_0 - (x^k+y^j)C_1 + x^k y^j C_2 = 0$ with Hankel--Wronskian square law $C_1^2 - C_0 C_2 = A B x^{m-2k}y^{n-2j}(x^k-y^j)^2$, the non-singular Pad\'e--Thue Wronskian tower $\det\mathcal{M}_N = -(3N+1)X Y C^{2N+1}$, the cross-floor Wronskian identity $\det\mathcal{W}_{2,3} = (U-V)^3(U^2+3UV+V^2)$ (\texttt{thm11}), and the root-free cubic Pad\'e compression identity $U I_1 \Phi_3(a\widehat{P}_1, b\widehat{Q}_1) = a^3(U-V)^3(U+V) - (U b^3 - V a^3)\widehat{Q}_1^3$ across all $d^2=9$ residue cells in $\mathbb{Z}$ (\texttt{thm10}). All theorems (\texttt{thm6}, \texttt{thm7}, \texttt{thm10}, \texttt{thm11}) are formally verified in pure \textbf{Lean~4} core (\texttt{0 sorryAx}, \texttt{0 Classical.choice}).

\medskip\noindent
\textbf{Keywords:} exponential Diophantine equations, Pillai's equation, Pad\'e--Thue approximation, Wronskian determinant, isotropic discriminant cone, parabolic $\mathrm{SL}_2(\mathbb{Z})$ Jordan matrices, Lean~4.

\smallskip\noindent
\textbf{MSC 2020:} 11D61, 11J68, 15B36, 68V20.
\end{abstract}

\section{Joint Matrix Decomposition on the Isotropic Discriminant Cone and the Parabolic \texorpdfstring{$\mathrm{SL}_2(\mathbb{Z})$}{SL2(Z)} Jordan Lock Quartet}
\label{sec:isotropic-cone}

\begin{theorem}[Joint Matrix Decomposition, Isotropic Discriminant Cone Generators, Exact Generator Factorization $C = (x^{m-a}\mp y^b)(x^a-y^b)$, and the Parabolic $\mathrm{SL}_2(\mathbb{Z})$ Jordan Lock Quartet $k\in\{6,12,15,40\}$]
\label{thm:main-6}
For the exponential Diophantine equation $A x^m - B y^n = C$ ($m,n\ge 2$), the following matrix and quadric-cone decompositions hold (\texttt{thm6\_pillai\_skeletal\_matrix\_decomposition\_and\_parabolic\_sl2\_jordan\_triptych}):
\begin{enumerate}[label=\textbf{(\roman*)},leftmargin=*]
\item \textbf{Joint matrix factorization and rank-$1$ cross-stage transition identity:}
  every interior exponent pair $(a,b)$ ($1\le a\le m-1$, $1\le b\le n-1$) defines a $2\times 2$ joint matrix $M_{a,b} = \begin{pmatrix} A x^{m-a} & B y^{n-b} \\ y^b & x^a \end{pmatrix}$ with invariant determinant $\det M_{a,b} = C$. The normalized cross-stage transition matrix $\widetilde{P} = \frac{1}{C} M_{a,b}^{\mathrm{adj}} M_{a',b'}\in\mathrm{SL}_2(\mathbb{Q})$ ($\det\widetilde{P} = +1$, $\Delta a = a'-a$, $\Delta b = b'-b$) obeys the exact rank-$1$ factorization
  \[
    \widetilde{P} \;-\; \begin{pmatrix} y^{\Delta b} & 0 \\ 0 & x^{\Delta a} \end{pmatrix}
    \;=\; \frac{x^{\Delta a}y^{\Delta b}-1}{C}\,\begin{pmatrix} -x^a \\ y^b \end{pmatrix}\begin{pmatrix} A x^{m-a'} & B y^{n-b'} \end{pmatrix}.
  \]
\item \textbf{Isotropic discriminant cone generators, exact generator factorization, and the $\mathrm{SL}_2(\mathbb{Z})$ nilpotent Jordan lock quartet ($k\in\{6,12,15,40\}$):}
  in the quadric coordinates $X = (\widetilde{P}_{11}-\widetilde{P}_{22})/2$, $Y = (\widetilde{P}_{12}+\widetilde{P}_{21})/2$, $T = (\widetilde{P}_{12}-\widetilde{P}_{21})/2$, the characteristic discriminant equals $\frac{1}{4}(\operatorname{Tr}\widetilde{P}^2 - 4) = X^2 + Y^2 - T^2$. On symmetric midpoint stages ($m=a+a'$, $n=b+b'$, $A=B=1$), $Y\equiv 0$ and the linear generator (ruling) $X - T$ obeys $2C(X - T) = (x^a - y^b)^2 - (x^{a'} - y^{b'})^2$, while $y^b(x^{m-a}-y^{n-b}) + x^{m-a}(x^a-y^b) = C$ yields the exact generator factorization $C = (x^{m-a}\mp y^b)(x^a - y^b)$ along the isotropic cone $X^2 - T^2 = 0$ ($\operatorname{Tr}\widetilde{P} = \mp 2$). For coprime bases, integrality $\widetilde{P}_k\in\mathrm{SL}_2(\mathbb{Z})$ generates the four canonical Jordan locks:
  \[
    \begin{aligned}
      3^2 - 2^3 &= 1 = 1\times 1 &\implies k &= 6 = |\square(3)|, &\quad \widetilde{P}_{\mathrm{I}} &= \begin{pmatrix} -7 & -6 \\ 6 & 5 \end{pmatrix} \quad (6^2 = 36,\ \operatorname{Tr}=-2), \\
      2^5 - 3^3 &= 5 = 5\times 1 &\implies k &= 12 = |V(O_2)|, &\quad \widetilde{P}_{\mathrm{II}} &= \begin{pmatrix} -13 & -12 \\ 12 & 11 \end{pmatrix} \quad (12^2 = 144,\ \operatorname{Tr}=-2), \\
      5^3 - 3^4 &= 44 = 22\times 2 &\implies k &= 15 = 3\cdot 5, &\quad \widetilde{P}_{\mathrm{III}} &= \begin{pmatrix} -16 & -15 \\ 15 & 14 \end{pmatrix} \quad (30^2 = 900,\ \operatorname{Tr}=-2), \\
      2^{10} - 5^4 &= 399 = 133\times 3 &\implies k &= 40 = 4^2+5^2-1, &\quad \widetilde{P}_{\mathrm{IV}} &= \begin{pmatrix} -39 & -40 \\ 40 & 41 \end{pmatrix} \quad (2^7-2^3 = 5^3-5^1 = 120,\ \operatorname{Tr}=+2).
    \end{aligned}
  \]
\end{enumerate}
\end{theorem}

\section{Three-Level Hankel--Wronskian Descent, the Pad\'e--Thue Wronskian Tower, and Root-Free Cubic Compression}
\label{sec:pade-thue}

\begin{theorem}[Unimodular Farey--B\'ezout exponent matrix, three-level recursive descent, Pad\'e--Thue gear-cascade monoid, adjacent-floor non-degeneracy $|I_{\mathcal{R}}| + |I_{\mathcal{R}\mathbin{+\!+}[r]}| \ge 1$, and bilateral Pad\'e bracket, \texttt{thm7}]
\label{thm:main-7}
For any coprime exponent pair $\gcd(m,n)=1$ of $A x^m - B y^n = C$, the Farey--B\'ezout parent stages $(a,b)$ and $(a',b') = (m-a,n-b)$ form $E_{\mathrm{B\acute{e}zout}} = \begin{pmatrix} a & b \\ a' & b' \end{pmatrix}\in\mathrm{SL}_2(\mathbb{Z})$ ($a b' - a' b = +1$), reconstructing the bases via $(x^a)^{b'} = x\cdot (x^{a'})^b$ and $(y^{b'})^a = y\cdot (y^b)^{a'}$. Along any ray $(k,j)$, consecutive stage differences $C_t = A x^{m-tk} - B y^{n-tj}$ satisfy the weight-independent three-term recurrence $C_0 - (x^k+y^j)C_1 + x^k y^j C_2 = 0$ and the Hankel--Wronskian square identity $C_1^2 - C_0 C_2 = A B x^{m-2k}y^{n-2j}(x^k-y^j)^2$. Across the countable monoid $\mathbb{Z}^{<\infty}$ of $M$-floor Pad\'e--Thue gear cascades $\mathcal{R} = [r_0,r_1,\dots,r_{M-1}]$ (\texttt{padeTowerP}, \texttt{padeTowerQ}, \texttt{padeTowerProdPlus}, \texttt{padeTowerProdMinus}), composition satisfies $P_{\mathcal{R}_1\mathbin{+\!+}\mathcal{R}_2}(C,B) = P_{\mathcal{R}_2}(P_{\mathcal{R}_1},Q_{\mathcal{R}_1})$, $Q_{\mathcal{R}_1\mathbin{+\!+}\mathcal{R}_2}(C,B) = Q_{\mathcal{R}_2}(P_{\mathcal{R}_1},Q_{\mathcal{R}_1})$, $\Pi_\pm(\mathcal{R}_1\mathbin{+\!+}\mathcal{R}_2) = \Pi_\pm(\mathcal{R}_1)\Pi_\pm(\mathcal{R}_2)$ ($\Pi_+(\mathcal{R})=\prod_{m=0}^{M-1}r_m$, $\Pi_-(\mathcal{R})=\prod_{m=0}^{M-1}(2-r_m)$), $P_{\mathcal{R}}+Q_{\mathcal{R}} = \Pi_+(\mathcal{R})(C+B)$, $P_{\mathcal{R}}-Q_{\mathcal{R}} = \Pi_-(\mathcal{R})A$ ($A = C - B$), and $Q_{\mathcal{R}}(u-v) - I_{\mathcal{R}} = -u\,\Pi_-(\mathcal{R})A$ for $I_{\mathcal{R}} := u P_{\mathcal{R}} - v Q_{\mathcal{R}}$. Moreover, between any floor $\mathcal{R}$ and the next floor $\mathcal{R}' = \mathcal{R}\mathbin{+\!+}[k+1]$ ($P_{\mathcal{R}'} = P_{\mathcal{R}} + k Q_{\mathcal{R}}$, $Q_{\mathcal{R}'} = k P_{\mathcal{R}} + Q_{\mathcal{R}}$):
\begin{equation}
  \label{eq:adjacent-pade-wronskian}
  \begin{aligned}
    I_{\mathcal{R}'} - (k+1)I_{\mathcal{R}} &= -k(u+v)\,\Pi_-(\mathcal{R})\,A, &\quad
    I_{\mathcal{R}'} + (k-1)I_{\mathcal{R}} &= k(u-v)\,\Pi_+(\mathcal{R})(C+B), \\
    Q_{\mathcal{R}'}I_{\mathcal{R}} - Q_{\mathcal{R}}I_{\mathcal{R}'} &= k\,u\,\Pi_-(\mathcal{R})\,\Pi_+(\mathcal{R})\,A(C+B), &\quad
    P_{\mathcal{R}'}Q_{\mathcal{R}} - Q_{\mathcal{R}'}P_{\mathcal{R}} &= -k\,\Pi_-(\mathcal{R})\,\Pi_+(\mathcal{R})\,A(C+B).
  \end{aligned}
\end{equation}
Consequently, whenever $k(u+v)\,\Pi_-(\mathcal{R})A \ne 0$, two adjacent-floor collision residuals never vanish simultaneously:
\[
  \neg\bigl(I_{\mathcal{R}} = 0 \;\land\; I_{\mathcal{R}'} = 0\bigr), \qquad
  |I_{\mathcal{R}}| + |I_{\mathcal{R}'}| \;\ge\; 1 \quad\bigl(\max(|I_{\mathcal{R}}|,|I_{\mathcal{R}'}|) \ge 1\bigr),
\]
and opposite signs between adjacent floors ($I_{\mathcal{R}} \le 0 \le I_{\mathcal{R}'}$ or $I_{\mathcal{R}'} \le 0 \le I_{\mathcal{R}}$) trap $(u,v)$ in the bilateral integer bracket $u P_{\mathcal{R}} \le v Q_{\mathcal{R}} \land v Q_{\mathcal{R}'} \le u P_{\mathcal{R}'}$ (respectively $u P_{\mathcal{R}'} \le v Q_{\mathcal{R}'} \land v Q_{\mathcal{R}} \le u P_{\mathcal{R}}$).
\end{theorem}

\begin{theorem}[Root-free $9$-cell cubic Pad\'e compression identity, \texttt{thm10}]
\label{thm:main-10}
For all $(U,V,a,b)\in\mathbb{Z}^4$, the collision invariant $I_1 = a\widehat{P}_1 - b\widehat{Q}_1$ ($\widehat{P}_1 = U+2V,\ \widehat{Q}_1 = 2U+V$) and the cubic norm factor $\Phi_3(x,y) = x^2 + xy + y^2$ satisfy
\begin{equation}
  \label{eq:root-free-cubic-compression}
  U\,I_1\,\Phi_3\bigl(a\widehat{P}_1,\,b\widehat{Q}_1\bigr) \;\equiv\; a^3(U-V)^3(U+V) \;-\; \bigl(U b^3 - V a^3\bigr)\,\widehat{Q}_1^3.
\end{equation}
\end{theorem}

\begin{theorem}[Cross-floor Wronskian identity between floors $k=2$ and $k=3$, \texttt{thm11}]
\label{thm:main-11}
For all $(U,V)\in\mathbb{Z}^2$, the cross-floor Wronskian between the quadratic Pad\'e pair ($\widehat{P}_{1,2}=U+3V,\ \widehat{Q}_{1,2}=3U+V$) and the cubic Pad\'e pair ($\widehat{P}_{1,3}=U+2V,\ \widehat{Q}_{1,3}=2U+V$) factors identically as
\begin{equation}
  \label{eq:cross-floor-wronskian}
  \det\mathcal{W}_{2,3} \;=\; \widehat{Q}_{1,2}^2\widehat{P}_{1,3}^3 - \widehat{P}_{1,2}^2\widehat{Q}_{1,3}^3 \;\equiv\; (U - V)^3\bigl(U^2 + 3UV + V^2\bigr).
\end{equation}
\end{theorem}

\section{Formal Verification in Lean~4}
\label{sec:lean4}

The formalization of Theorems~\ref{thm:main-6} and~\ref{thm:main-7} (\texttt{thm6}, \texttt{thm7}, \texttt{thm10}, \texttt{thm11}) resides in \texttt{PillaiDiophantine.lean} in \url{https://github.com/AI-Guiders/scientific-papers} (\texttt{0 sorryAx}, \texttt{0 Classical.choice}).

\begin{thebibliography}{9}
\bibitem{Pillai1931}
S.\,S.~Pillai, ``On the inequality $0 < a^x - b^y \le n$,'' \emph{Journal of the Indian Mathematical Society} \textbf{19} (1931), 1--11.
\bibitem{Thue1909}
A.~Thue, ``\"Uber Ann\"aherungswerte algebraischer Zahlen,'' \emph{Journal f\"ur die reine und angewandte Mathematik} \textbf{135} (1909), 284--305.
\bibitem{Moura2021}
L.~de~Moura and S.~Ullrich, ``The Lean~4 theorem prover and programming language,'' in \emph{CADE~28}, LNCS vol.~12699, Springer, 2021, pp.~625--635.
\end{thebibliography}

\end{document}
